\section{Keeping canonical cancellation inside an exact cell norm (NS-67)}
\label{sec:ns67-canonical-cells}

\paragraph{Status.}
The canonical growth criterion in Section~\ref{sec:ns66-canonical-growth}
is still open. Here the complete once-smoothed error is evaluated before
any absolute-value estimate is taken. This gives exact arithmetic cell
identities, a smaller elementary tail bound, and a finite growing-interval
version of the same growth target. It does not provide the missing
subpolynomial estimate. No finite table proves a cofinal bound.

Set $e_N=J(\chi+B_N)$ and $F_N(t)=e_N(1/t)$, $t>0$.
Write
\[
 b_N=Ns_N,\qquad c_{N,n}=\mu(n)-b_N\mathbf1_{n=N},\qquad
 q_N(j)=\sum_{\substack{n\mid j\\n\leq N}}\mu(n).
\]
The letter $b_N$ in this section denotes the canonical boundary
coefficient, not a target-load vector. The identity
$\sum_{n\leq N}c_{N,n}/n=0$ is exact.

\begin{proposition}[Complete physical cells]
\label{prop:ns67-cells}
We have $F_N(t)=0$ for $0<t\leq N$, and, for every integer $j\geq N$,
\begin{equation}
 F_N(t)=f_{N,j}+a_{N,j}\log(t/j),\quad j\leq t\leq j+1,
 \qquad
 a_{N,j}=1-\sum_{n\leq N}\mu(n)\left\lfloor\frac jn\right\rfloor
                   +b_N\left\lfloor\frac jN\right\rfloor.
 \label{eq:ns67-cell}
\end{equation}
The endpoint values are generated by
\[
 f_{N,N}=0,\qquad
 f_{N,j+1}=f_{N,j}+a_{N,j}\log(1+1/j).
\]
Equivalently, putting $d_N(j)=q_N(j)-b_N\mathbf1_{N\mid j}$,
\[
 a_{N,j}-a_{N,j-1}=-d_N(j),\qquad j>N.
\]
The initial slope is $a_{N,N}=b_N$, and therefore
\begin{equation}
 F_N(t)=b_N\log(t/N)
          -\sum_{N<m\leq t}d_N(m)\log(t/m),\qquad t\geq N.
 \label{eq:ns67-divisor-response}
\end{equation}
For $h_j=\log(1+1/j)$ and
$I_k(h)=\int_0^h u^ke^{-u}\,du$, the complete cell energy is
\begin{equation}
 C_{N,j}=\int_j^{j+1}\frac{F_N(t)^2}{t^2}\,dt
 =\frac{f_{N,j}^2 I_0(h_j)+2f_{N,j}a_{N,j}I_1(h_j)
                       +a_{N,j}^2 I_2(h_j)}j.
 \label{eq:ns67-cell-energy}
\end{equation}
In particular $D_{N,1}^2=\sum_{j=N}^\infty C_{N,j}$.
Every signed divisor contribution and every cross term remains present.
\end{proposition}
\begin{proof}
The support assertion is the exact exterior interpolation in NS-62.
For $t\geq1$ the complete expression is
\[
 F_N(t)=\log t+\sum_{n\leq N}c_{N,n}A(t/n),
 \qquad A(z)=z-k\log z+\log(k!),\quad k=\lfloor z\rfloor.
\]
On a unit cell all floors are constant. The coefficient of $t$ is
$\sum c_{N,n}/n=0$, and the coefficient of $\log t$ is precisely
$a_{N,j}$. Continuity of $A$ at integers gives the endpoint recurrence.
The floor jumps give the divisor expression; at the initial endpoint
$\sum_{n\leq N}\mu(n)\lfloor N/n\rfloor=1$, so $a_{N,N}=b_N$.
The change $t=je^u$ proves~\eqref{eq:ns67-cell-energy}.
\end{proof}

\begin{lemma}[An elementary complete Stirling remainder]
\label{lem:ns67-remainder}
For every real $z\geq1$,
\begin{equation}
 \left|A(z)-\tfrac12\log(2\pi z)\right|\leq\frac1{6z}.
 \label{eq:ns67-remainder}
\end{equation}
\end{lemma}
\begin{proof}
Let $R(z)=A(z)-\tfrac12\log(2\pi z)$ and
$\widetilde B_2(z)=\{z\}^2-\{z\}+1/6$.
Then $R'(z)=(\{z\}-1/2)/z$ almost everywhere.
Stirling's formula gives $R(z)\to0$ as $z\to\infty$.
Since $\widetilde B_2$ is continuous at integers and its derivative
is $2(\{z\}-1/2)$ off the integers, integration by parts gives
\[
 R(z)=\frac{\widetilde B_2(z)}{2z}
             -\frac12\int_z^\infty\frac{\widetilde B_2(u)}{u^2}\,du.
\]
Both terms are controlled by $|\widetilde B_2|\leq1/6$, proving
\eqref{eq:ns67-remainder}. There is no unestimated periodic tail.
\end{proof}

Define the exact signed leading coefficients and the absolute remainder
mass by
\[
 A_N=1+\tfrac12\sum_{n\leq N}c_{N,n}=1-\tfrac12N\eta_N,
 \qquad B_N^\flat=\tfrac12\sum_{n\leq N}c_{N,n}\log\frac{2\pi}{n},
 \qquad L_N=\sum_{n\leq N}n|c_{N,n}|.
\]
The superscript distinguishes $B_N^\flat$ from the canonical function $B_N$.
Then for $t\geq N$,
\begin{equation}
 F_N(t)=A_N\log t+B_N^\flat+R_N^\flat(t),\qquad
 |R_N^\flat(t)|\leq\frac{L_N}{6t}.
 \label{eq:ns67-full-tail}
\end{equation}
For $T\geq N$ put $v_N(T)=A_N\log T+B_N^\flat$ and
\[
 P_N(T)=\frac{v_N(T)^2+2A_Nv_N(T)+2A_N^2}{T},\qquad
 r_N(T)=\frac{L_N}{6\sqrt{3T^3}}.
\]
The entire tail obeys
\begin{equation}
 \left|\int_T^\infty\frac{F_N(t)^2}{t^2}\,dt-P_N(T)\right|
       \leq2\sqrt{P_N(T)}\,r_N(T)+r_N(T)^2.
 \label{eq:ns67-tail-energy}
\end{equation}
This follows by integrating the squared linear logarithm exactly and
applying Cauchy--Schwarz to its cross term with the bounded remainder.
The right side includes the remainder square; no cancellation with it
is presumed.

\begin{proposition}[The polynomial cutoff retains the growth target]
\label{prop:ns67-polynomial-cutoff}
Let
\[
 Q_N=\left(\sum_{j=N}^{N^2-1}C_{N,j}\right)^{1/2},\qquad N\geq2.
\]
Unconditionally,
\begin{equation}
 0\leq D_{N,1}-Q_N\ll\log(2N).
 \label{eq:ns67-cutoff-error}
\end{equation}
Consequently
\[
 \lim_{N\to\infty}\frac{\log(1+Q_N)}{\log N}=\beta_*-\frac12,
\]
and RH is equivalent to $Q_N=O_\varepsilon(N^\varepsilon)$ for every
$\varepsilon>0$. This is an exact finite growing-interval reformulation,
not a proven bound for $Q_N$.
\end{proposition}
\begin{proof}
The elementary divisor identity gives
\[
 b_N=1+\sum_{n\leq N}\mu(n)\{N/n\},\qquad |b_N|\leq N+1.
\]
Hence $\sum|c_{N,n}|\leq2N+1$, $|A_N|\leq N+3/2$,
$|B_N^\flat|\ll N\log(2N)$, and $L_N\ll N^2$.
At $T=N^2$, therefore, $P_N(T)^{1/2}=O(\log(2N))$ and
$r_N(T)=O(1/N)$. Equation~\eqref{eq:ns67-tail-energy} bounds the
complete tail norm by $O(\log(2N))$. Disjoint-support orthogonality
gives $D_{N,1}^2=Q_N^2+\text{tail energy}$ and thus
\eqref{eq:ns67-cutoff-error}.
Apply the exponent theorem~\ref{cor:ns66-growth-exponent}. When
$\beta_*=1/2$, $0\leq Q_N\leq D_{N,1}$ gives exponent zero.
When $\beta_*>1/2$, subtracting an $O(\log N)$ error does not change
any positive power lower bound coming from a zero to the right of the
line. Take the supremum over those zeros and combine with the upper
bound $Q_N\leq D_{N,1}$.
\end{proof}

\paragraph{Complete boundary/integral comparison.}
Write $w_N=N\eta_N$ and
$I_N=\int_1^N M(u)\rho_u\,du/u$.
Then the exact Abel identity is $e_N=\tau+I_N-w_N\sigma_N$.
Once the full energy and the mixed integrals have been enclosed,
\[
 \|\tau+I_N\|^2=D_{N,1}^2+2w_N\langle e_N,\sigma_N\rangle
                          +\frac{w_N^2}{N}\|\sigma_1\|^2,
\]
\[
 \|I_N\|^2=\|\tau+I_N\|^2+2
       -2\{\langle e_N,\tau\rangle+w_N\langle\sigma_N,\tau\rangle\}.
\]
These are comparisons of the complete terms, not of their absolute
integrands. The implementation integrates both mixed terms on every
cell and retains their complete tails. The identities make the
cancellation measurable without discarding it; they do not bound it
uniformly in $N$.

\begin{proposition}[The two complete Abel terms have a uniform conditioning bound]
\label{prop:ns67-abel-conditioning}
Put $c_\sigma=\|\sigma_1\|$ and $k_N=k_{N,1}$ from
Proposition~\ref{prop:ns62-smoothed-obstruction}. With
$U_N=\tau+I_N$ and $V_N=w_N\sigma_N$, we have
\begin{equation}
 \|V_N\|\leq c_\sigma D_{N,1}+\frac{2c_\sigma}{\sqrt N},\qquad
 D_{N,1}\leq\|U_N\|+\|V_N\|
 \leq(1+2c_\sigma)D_{N,1}+\frac{4c_\sigma}{\sqrt N}.
 \label{eq:ns67-conditioning}
\end{equation}
Moreover define $\widetilde V_N=(w_N-k_N)\sigma_N$ and
$\widetilde U_N=U_N-k_N\sigma_N$. Then
$e_N=\widetilde U_N-\widetilde V_N$ and
\begin{equation}
 \|\widetilde V_N\|\leq c_\sigma D_{N,1},\qquad
 D_{N,1}\leq\|\widetilde U_N\|+\|\widetilde V_N\|
                \leq(1+2c_\sigma)D_{N,1}.
 \label{eq:ns67-centered-conditioning}
\end{equation}
Thus cancellation between these two complete vectors cannot change the
norm-growth exponent. Their internal arithmetic cancellation remains
unestimated. This does not justify replacing the integral $I_N$ by
its absolute integrand.
\end{proposition}
\begin{proof}
The localized unitary witness already proved in NS-62 gives
$|w_N-k_N|/\sqrt N\leq D_{N,1}$.
The defining integral of $k_N$ has absolute value at most $2$, since
$|\sin(2\pi x)/(\pi x)|\leq2$ and
$N\int_0^{1/N}\log(1/(Nx))\,dx=1$.
Use $\|\sigma_N\|=c_\sigma/\sqrt N$ and the triangle inequality.
The centered inequalities follow with no additive error; both centered
vectors differ by exactly $e_N$, so no cross term is omitted.
\end{proof}

\paragraph{Complete finite checks.}
The physical cell implementation evaluates $N=8,16,64,512,4096$ at
256 and 384 bits, with the entire tail beyond $T=1024N$ enclosed by
\eqref{eq:ns67-tail-energy}. It also retains the mixed tails with
$\sigma_N$ and $\tau$. Independent complete Gram evaluations at
$N=8,16,64$ lie inside all three corresponding physical enclosures.
Doubling the physical cutoff at $N=16,512$ gives overlapping full
energies and mixed terms. The replay verifies $175$ cross-precision
quantities, $28$ cutoff comparisons and $9$ independent Gram comparisons.
In all five cases the centered ratio in
\eqref{eq:ns67-centered-conditioning} lies below the certified
constant $1+2c_\sigma<4.617$. These are complete finite norms;
no bounded cofinal subsequence is inferred.

\paragraph{What changed in the proposed cancellation route.}
The large finite ratios from separately bounding the two original Abel
vectors can include a harmless $N^{-1/2}$ component. Removing its explicit
coefficient $k_N$ makes their comparison uniformly controlled by the
actual error. Therefore the proposed search for a power-saving hidden
solely in cancellation between those two \emph{complete} vectors does
not supply a weaker arithmetic target. The possible useful cancellation
is inside the integral, or inside the exact divisor-cell response
\eqref{eq:ns67-divisor-response}. No bound for either signed quantity
is proved here. This is an obstruction to that specific decomposition
argument, not a failure of the canonical subpolynomial-growth target.
